\section{Responsabilidades en el ecosistema: plataformas, proveedores de tecnología, organismos de denuncia e investigadores}

Por último, volvemos al sistema de extremo a extremo. En el ecosistema de seguridad infantil ningún participante tiene todas las capacidades: la plataforma controla la relación con los usuarios y las action del producto; los proveedores de tecnología aportan las detection/review primitives; los organismos de denuncia previstos por la ley reciben las pistas y las complementan con más información; las autoridades y los servicios de atención a víctimas investigan y protegen conforme a la ley. Tener responsabilidades claras no es eludirlas. Sirve para evitar que el proveedor del modelo se exceda en sus atribuciones, que la plataforma subcontrate su gobernanza o que los investigadores se vean rebasados por datos de baja calidad.

\subsection{Límites de responsabilidad y ciclo de retroalimentación}

La plataforma debe hacerse cargo de la policy, las notificaciones a usuarios, las apelaciones, las medidas sobre cuentas y las obligaciones legales. Los proveedores de tecnología como Thorn deben ofrecer herramientas validadas y auditables, y proteger los datos sensibles. Los organismos de denuncia como el NCMEC procesan los reportes y las pistas según el papel que les asigna la ley. Las autoridades se encargan de la investigación, las pruebas y los procedimientos legales. Las organizaciones de apoyo a víctimas brindan protección continua. Un verified outcome solo puede retroalimentar al modelo y al reference system si hay autorización, minimización de datos y limitación de la finalidad.

\lecturefigure{16-ecosystem-roles.png}{Los resultados en seguridad infantil dependen de una división clara del trabajo entre plataformas, proveedores de tecnología, organismos de denuncia e investigadores.}{Subtítulos locales 37:10--37:53; reconstrucción conceptual basada en la closing synthesis de la clase.}

\begin{knowledgebox}{Lectura de la figura: cada handoff necesita un contract}
El contract entre platform y provider establece las entradas, las salidas, el SLA, el uso de los datos y los incident; el contract entre platform y reporting agency establece los campos que exige la ley, el manejo de duplicados y la cadena de custodia de las pruebas; el handoff de la agency a los investigators requiere prioritization y jurisdiction; cuando el investigation outcome regresa, su uso debe quedar restringido. Un handoff ambiguo genera reportes duplicados, campos faltantes, casos que nadie atiende o accesos que exceden las atribuciones.
\end{knowledgebox}

\begin{importantbox}{Implicaciones de ingeniería de un mission-driven business model}
Una misión sin fines de lucro y el software comercial no entran en conflicto por naturaleza: los productos de pago pueden financiar la investigación y el desarrollo continuos, el soporte, las auditorías y una infraestructura de alta disponibilidad, mientras que la gobernanza basada en la misión limita la elección de clientes, el uso de los datos y las prioridades del producto. El criterio real de aceptación no es la etiqueta de la organización, sino si puede mantener herramientas de seguridad a largo plazo, hacer públicos sus límites, aceptar supervisión y hacer que la métrica principal sea el resultado en protección y no el engagement.
\end{importantbox}

\teachervoice{El cierre de Cordua describe los papeles sin rodeos: la plataforma debe detectar y tomar medidas, y los investigadores deben identificar a las víctimas y dar seguimiento a los casos. La función de la tecnología es que la información que pasa de unos a otros sea más rápida, menos repetida y más útil para actuar, no que un modelo sustituya a las instituciones y a los profesionales.}

\subsection{Resumen de la sección}

La seguridad infantil es un sistema que abarca varias organizaciones. Con roles, data contracts, audit trails y límites de retroalimentación claros, cada parte puede aprovechar su especialidad. Si las responsabilidades son ambiguas, el riesgo se traslada al eslabón que tiene menos contexto o menos recursos.
